Source : OpenFisca-France 176.1 (aujourd'hui, droit des revenus 2025) ; calculs des auteurs (note8p/calculs/ircsg\_*.py).\par Salarié vivant seul, sans enfant : prélèvement sur 100 € de revenu après cotisations en plus. Fusion naïve : barème actuel de l'IR (sans décote ni abattement) appliqué au revenu après cotisations, plus un taux proportionnel qui donne la même recette. La CEHR (3 \%, puis 4 \% au-delà de 500 000 € de revenu fiscal) s'ajoute aujourd'hui au-delà d'environ 270 000 €. Courbe d'aujourd'hui : médiane glissante sur ± 1 000 €, qui neutralise les arrondis à l'euro du calcul de l'IR.
